\chapter{INTRODUCCIÓN}
En la era digital actual, los jóvenes enfrentan una paradoja significativa: aunque están hiperconectados y acumulan numerosos contactos digitales, experimentan una creciente desconexión emocional. Las plataformas sociales contemporáneas priorizan métricas cuantitativas y la retención de usuarios, descuidando la calidad de las interacciones y el estado emocional de sus usuarios. Esta dinámica genera entornos donde las conexiones son abundantes pero superficiales, limitando las oportunidades para establecer vínculos auténticos.
La Interacción Humano-Computadora Afectiva (IHC Afectiva) busca transformar esta realidad mediante el diseño de interfaces que consideren los estados emocionales de los usuarios. Estudios demuestran que las aplicaciones diseñadas bajo estos principios, especialmente aquellas que incorporan micro interacciones positivas diarias, pueden generar efectos significativos en el estado emocional juvenil.
Este proyecto propone el desarrollo de una aplicación móvil fundamentada en principios de IHC Afectiva para fomentar micro interacciones sociales positivas entre jóvenes a través de mensajes breves diarios. La aplicación emparejará usuarios para intercambios limitados pero significativos, e incorporará avatares personalizables para expresar identidad y emociones. Esta herramienta busca contribuir a la regulación emocional de los jóvenes a través de micro interacciones positivas diarias, favoreciendo progresivamente la mejora de su estado emocional en entornos digitales, fortalecer vínculos genuinos y promover comunidades más empáticas, priorizando la calidad de la interacción sobre la cantidad.

\section{ANTECEDENTES}
La Interacción Humano-Computadora Afectiva (IHC Afectiva) es una subárea de la IHC que se enfoca en el diseño de interfaces que consideren y respondan a los estados emocionales de los usuarios, con el objetivo de crear experiencias más empáticas y significativas.
En investigaciones realizadas a nivel internacional, se encontró a Baudon y Jachens (2021) quienes en su investigación desarrollada en Reino Unido, buscaron examinar el impacto de aplicaciones móviles que promueven la regulación emocional, salud mental positiva y bienestar en la población general. Para esto realizaron una revisión sistemática y meta-análisis de 52 estudios que describían 48 aplicaciones, con un total de 22,090 participantes, aplicando criterios de inclusión para aplicaciones que promovieran emociones positivas, considerando aspectos como: regulación emocional, promoción de bienestar, efectos en salud mental positiva y usabilidad de las aplicaciones. En esta investigación, se presentó que las aplicaciones móviles de salud mental demostraron efectos pequeños pero significativos en la mejora del bienestar emocional, siendo particularmente efectivas aquellas que incorporaban elementos de micro-interacciones positivas diarias y técnicas de regulación emocional instantánea. De ahí que se concluye que las aplicaciones móviles pueden ser herramientas efectivas para promover emociones positivas y bienestar cuando están diseñadas con enfoques específicos de IHC Afectiva que consideran las micro-interacciones emocionales.

De igual forma, se encontró a Vakili (2024), quien en su estudio sobre preferencias de mensajería de texto para jóvenes en transición a servicios de salud mental, aplicó los modelos IMB y Diseño de Sistemas Persuasivos (PSD). Para ello, estudió a 100 jóvenes canadienses de entre 16 y 26 años, aplicando una encuesta que consideró aspectos como el contenido de los mensajes, las características tecnológicas y los facilitadores del \textit{engagement}.

En esta investigación se evidenció que la personalización de los mensajes fue el factor más significativo para el \textit{engagement} (27\%), seguido por los mensajes que captan la atención (20\%) y un número razonable de mensajes (15\%). Asimismo, se observó un mayor consenso en los elementos conductuales (100\%) que en los informativos (44\%) o motivacionales (25\%).

Se concluye que el diseño centrado en el usuario es un elemento fundamental en aplicaciones de Interacción Humano-Computadora (IHC) Afectiva que buscan generar microinteracciones significativas.

A nivel conceptual, se encontró el estudio "Social connection as a key target for youth mental health" (2025) el cual buscó establecer las conexiones sociales positivas como objetivo clave para la salud mental juvenil. Para esto se realizó una revisión de evidencia sobre conexiones entre pares y conexión escolar, considerando aspectos como: intervenciones de redes sociales, programas escolares y factores transdiagnósticos. En esta investigación, se presentó que las intervenciones de redes sociales y la conexión escolar muestran promesa para mejorar la salud mental juvenil, actuando las conexiones sociales como factor transdiagnóstico. De ahí que se concluye que las estrategias de prevención deberían enfocarse en la conexión social juvenil como elemento fundamental para el estado emocional.

\section{DESCRIPCIÓN DEL PROBLEMA}
En la actualidad, los jóvenes se encuentran en una paradoja digital: están hiperconectados a través de redes sociales, acumulando decenas o cientos de contactos, pero experimentan una desconexión emocional y social significativa debido a interacciones superficiales que priorizan métricas cuantitativas sobre conexiones genuinas.Esta situación refleja una creciente dificultad para establecer vínculos significativos en entornos digitales, donde la conexión constante no necesariamente se traduce en vínculos significativos.
Las plataformas digitales actuales se enfocan en la visibilidad masiva y la retención de usuarios, descuidando el estado emocional y generando entornos que carecen de experiencias significativas para los jóvenes. Este fenómeno se origina en el diseño de las plataformas sociales existentes, las cuales están orientadas a maximizar el tiempo de uso mediante algoritmos que promueven el consumo pasivo de contenido, en lugar de fomentar vínculos auténticos entre los usuarios. Como resultado, se produce una brecha creciente entre la cantidad de conexiones digitales y la calidad de las relaciones interpersonales que los jóvenes logran establecer.
Por otra parte, existe una ausencia de herramientas digitales que consideren los estados emocionales de los usuarios y sus necesidades de conexión social genuina, lo cual repercute en la falta de espacios que prioricen la calidad sobre la cantidad en las interacciones. Además, las dinámicas de comunicación actuales favorecen la unidireccionalidad y las respuestas automáticas, limitando las oportunidades para generar diálogos significativos que contribuyan al estado emocional juvenil.
En consecuencia, los jóvenes experimentan un deterioro progresivo en su estado emocional, manifestado en sentimientos de soledad, ansiedad social y dificultades para establecer vínculos auténticos tanto en el entorno digital como presencial. Esta situación perpetúa las interacciones superficiales como norma en la comunicación juvenil y dificulta el desarrollo de comunidades digitales empáticas y solidarias.

\subsection{Definición del problema}
Formulación del problema
¿Cómo abordar eficazmente la desconexión emocional que experimentan los jóvenes en la era digital al interactuar con las plataformas sociales existentes?

\section{OBJETIVOS DEL PROYECTO}
A continuación, se presentan el objetivo general y los objetivos específicos.

\subsection{Objetivo General}
Diseñar e implementar una aplicación móvil basada en principios de Interacción Humano-Computadora Afectiva que promueva la generación de micro interacciones sociales positivas entre jóvenes, contribuyendo a mejorar el estado emocional en entornos digitales.

\subsection{Objetivo Específicos}
\begin{itemize}
    \item Analizar el nivel de desconexión emocional que experimentan los jóvenes al interactuar con plataformas digitales, identificando factores que limitan la generación de vínculos significativos.
    \item Diseñar una interfaz móvil centrada en el usuario que incorpore dinámicas de interacción breves y elementos visuales orientados a fortalecer la conexión emocional.
    \item Implementar funcionalidades que permitan la interacción diaria entre los usuarios, como el envío y recepción de mensajes, junto con la personalización del avatar.
    \item Realizar pruebas piloto con un grupo de jóvenes para evaluar la usabilidad, la experiencia de usuario (UX) y la influencia de la aplicación en el estado emocional de los usuarios.
\end{itemize}

\section{JUSTIFICACIÓN}
El estado emocional de los jóvenes constituye una preocupación creciente en la era digital, donde la hiperconectividad no siempre se traduce en experiencias afectivas significativas. Aunque las redes sociales facilitan la comunicación masiva, con frecuencia generan una desconexión emocional, en la que la cantidad de contactos no asegura vínculos auténticos ni relaciones de apoyo genuinas. Esta situación evidencia la necesidad de herramientas digitales que promuevan la calidad de las interacciones y favorezcan la conexión emocional entre los usuarios.
En este contexto, el desarrollo de una aplicación móvil orientada a fomentar micro interacciones sociales positivas busca atender esta necesidad, promoviendo conexiones emocionales reales entre los usuarios. La plataforma se centrará en reforzar la empatía, la reciprocidad y la comunicación genuina, permitiendo que los jóvenes experimenten interacciones significativas que favorezcan la regulación emocional en el momento y, de forma progresiva, impacten positivamente en su estado emocional. Además, esta aplicación busca ofrecer un entorno seguro y positivo, donde las interacciones sean pequeñas pero significativas. 
Este proyecto surge como respuesta a la falta de herramientas digitales que prioricen la calidad de las interacciones sobre la cantidad. Las plataformas tradicionales suelen enfocarse en la visibilidad masiva y la retención de usuarios, lo que puede generar entornos superficiales que no consideran las emociones ni el impacto en el estado emocional de los jóvenes. 
Con este desarrollo, se espera reducir la brecha entre la hiperconectividad digital y el estado emocional de los jóvenes, promover la creación de espacios digitales más empáticos y fomentar interacciones auténticas y significativas que contribuyan a su estado emocional.

\section{ALCANCE Y LÍMITES}
Alcances:

\begin{itemize}
    \item El proyecto desarrollará un prototipo de una aplicación móvil funcional e interactiva centrada en fomentar micro interacciones sociales positivas entre jóvenes.
    \item Se incluirán funcionalidades de envío y recepción de mensajes breves diarios diseñados para fomentar la conexión emocional.
    \item Se implementará un sistema de personalización de avatares que permita la expresión identitaria de los usuarios.
    \item Se diseñará la interfaz centrada en el usuario con elementos visuales orientados a fortalecer la conexión emocional.
    \item Se realizarán pruebas piloto con un grupo de jóvenes de entre 16 y 28 años para evaluar la usabilidad, la experiencia de usuario (UX) y el impacto de la aplicación en la regulación emocional de los usuarios.
\end{itemize}

Límites:

\begin{itemize}
    \item La aplicación no sustituirá intervenciones terapéuticas profesionales ni funcionará como tratamiento para trastornos de salud mental.
    \item No se incluirán métodos clínicos formales ni pruebas que midan directamente mejoras en la salud mental de los usuarios.
    \item No se abordarán funcionalidades de redes sociales masivas como transmisiones en vivo, publicaciones públicas o métricas de popularidad.
    \item El enfoque se centrará únicamente en micro interacciones entre usuarios; no se desarrollarán recursos para comunicación grupal o individual masiva.
    \item La funcionalidad estará limitada a la plataforma móvil específica de Android; no se desarrollará versión web en esta primera fase.
\end{itemize}

\section{DISEÑO METODOLÓGICO}
FALTA COPIAR DE DOCS GOOGLE

